\section{Sublevels in spherical residues}
\label{sec:residues}

Inside a spherical coset, we consider smaller cosets whose exposed layer
counts lie below the ambient harmonic vector.  The two-point estimate will
show that a nonempty sublevel has only one source in its vertex filtration.
The isolated-layer obstruction has a different role: for the closed
inequalities associated with a positive null vector, it guarantees
nonemptiness and excludes the rank-zero exterior data.  We begin with the
topological principle that converts the link calculations into
contractibility.

\begin{lemma}[Vertex filtration]
\label{lem:filtration-principle}
Let $\mathcal P$ be a set-sized poset, and well-order its elements by
$\prec$.  For $v\in\mathcal P$, define the earlier lower and upper posets
\[
 \mathcal L_v=\{w\in\mathcal P:w<v,\ w\prec v\},\qquad
 \mathcal U_v=\{w\in\mathcal P:v<w,\ w\prec v\}.
\]
The link along which $v$ is added to the induced order complex on its
predecessors is the simplicial join of the order complexes of
$\mathcal L_v$ and $\mathcal U_v$.  If one poset is empty, this link is
the order complex of the other; if both are empty, the link is empty.

Suppose each attaching link is empty or contractible.  Then every
component of $|\mathcal P|$ is contractible and contains exactly one
\emph{source}, meaning a vertex with empty attaching link.  In
particular, the conclusion holds if each of $\mathcal L_v$ and
$\mathcal U_v$ is empty or contractible for every $v$.
\end{lemma}

\begin{proof}
A chain containing $v$ consists of a chain below $v$, the vertex $v$,
and a chain above $v$.  Every element of the first chain precedes every
element of the second in the poset order.  Restricting all other vertices
to those earlier than $v$ proves the asserted join description.  The
join of two nonempty contractible complexes is contractible.  Together
with the stated conventions for empty factors, this proves the last
assertion from the first one.

Index the vertices as $v_\alpha$, $\alpha<\kappa$, and let $K_\alpha$
be the induced order complex on $\{v_\beta:\beta<\alpha\}$.  At a
successor stage we attach the cone on the earlier link $L_\alpha$:
\[
 K_{\alpha+1}=K_\alpha\cup_{L_\alpha}(v_\alpha*L_\alpha).
\]
An empty link adds an isolated vertex and starts one component.  A
nonempty contractible link lies in a single component of $K_\alpha$.
The inclusion of that link in its cone is a homotopy equivalence and
a subcomplex cofibration.  The homotopy extension property therefore
makes the induced inclusion of the old component into the new one a
homotopy equivalence.  No other component changes, and two old
components cannot merge.  Thus the asserted description of components
and their sources persists at successor stages.

At a limit ordinal $\lambda$, the complex is the union of the earlier
subcomplexes, with the CW topology.  We recall why compact images cause
no difficulty in this union.  A compact subset of a closure-finite CW
complex lies in a finite subcomplex.  Otherwise it meets infinitely many
open cells; choose one point in each of a countable set of these cells.
Every subset of the chosen points meets each closed cell in a finite
set, by closure-finiteness, and hence is closed in the whole complex by
the weak topology.  The chosen points would therefore form an infinite
closed discrete subspace of the compact subset, a contradiction.  The
closures of the finitely many cells met by a compact subset form the
required finite subcomplex.

Consequently a path or a sphere map in $K_\lambda$ has image in some
$K_\alpha$ with $\alpha<\lambda$: only finitely many vertex indices
occur in its finite carrier.  Paths cannot merge previously distinct
components at the limit.  Every sphere map within a component is
null-homotopic in an earlier component, which is contractible by
induction.  Each component of $K_\lambda$ is thus weakly contractible
and has the same unique source as any of its earlier vertices.
Components are CW subcomplexes, so the CW Whitehead theorem makes them
contractible.  Transfinite induction, including the final union,
proves the result.
\end{proof}

We will also use the order homotopy observation that an order map
$f:\mathcal P\to\mathcal P$ satisfying $f(x)\le x$ for every $x$, or
$f(x)\ge x$ for every $x$, induces a map homotopic to the identity.
For example, when $f(x)\le x$, the assignments
$(x,0)\mapsto f(x)$ and $(x,1)\mapsto x$ form an order map from
$\mathcal P\times\{0<1\}$: if $x\le y$, then $f(x)\le y$.
The usual prism triangulation on every chain gives the homotopy, and
these homotopies agree on faces.  Thus an order retraction pointwise
comparable with the identity is a homotopy equivalence on order
complexes.  A poset with a maximum has a conical order complex.

\subsection{The two local sublevels}

Fix a spherical subset $T\subseteq S$, an integer $q$, and a right
coset $A_Tx$.  Assume that the layer data are capped at $q$ on
$\partial T$.  Operations in $A_T$ leave all exterior data unchanged,
so this assumption is independent of the representative.  Let $b$ be
the full vector obtained from $m^q(x)$ by harmonic extension on $T$:
\[
 b_i=m_i^q(x)\quad(i\notin T),\qquad (Cb)_T=0.
\]
The vector $b$ is nonnegative and independent of $x$ within its coset.
Here and throughout, caps concern all levels greater than $q$, whereas
the displayed inequalities concern level $q$ alone.

Consider pairs
\begin{equation}
 \label{eq:residues-pairs}
 X=(J,A_Jz),\qquad J\subsetneq T,\qquad z\in A_Tx,
\end{equation}
ordered by nested types and actual inclusion of cosets.  The open
colors of $X$ are $V=T\setminus J$.  Counts and caps on $V$ are
well-defined on $A_Jz$, since operations in $A_J$ change data only
in colors of $J$.

The \emph{strict sublevel} $\mathcal R_{<}(T,A_Tx;q)$ consists of
the pairs in \eqref{eq:residues-pairs} for which every $i\in V$ is
capped at $q$ and satisfies
\begin{equation}
 \label{eq:residues-strict}
 m_i^q(z)<b_i.
\end{equation}
We use a closed sublevel only with the following additional data:
\begin{enumerate}[label=(\roman*)]
 \item $q>0$;
 \item $I\subseteq S$ is nonempty and connected, with $T\subseteq I$;
 \item $c_I>0$ and $C_Ic_I=0$, and $c$ is extended by zero outside $I$;
 \item on $(I\setminus T)\cup\partial I$ the data are capped at $q$
 and satisfy $m^q=c$.
\end{enumerate}
These are exterior conditions for $A_Tx$.  In particular they imply
the assumed cap on $\partial T$.  Harmonicity and uniqueness of
extension give $b_T=c_T$.  The \emph{special closed sublevel}
$\mathcal R_{\le}(T,A_Tx;q,I,c)$ consists of the pairs in
\eqref{eq:residues-pairs} for which every $i\in V$ is capped at $q$
and satisfies
\begin{equation}
 \label{eq:residues-closed}
 m_i^q(z)\le b_i=c_i.
\end{equation}
No connectedness assumption is imposed on $T$.

\begin{proposition}[Spherical-residue sublevels]
\label{prop:residue-sublevels}
For every nonempty spherical $T$ and the data just specified:
\begin{enumerate}[label=(\roman*)]
 \item $\mathcal R_{<}(T,A_Tx;q)$ is empty or contractible;
 \item whenever the additional exterior hypotheses hold,
 $\mathcal R_{\le}(T,A_Tx;q,I,c)$ is nonempty and contractible.
\end{enumerate}
For $T=\varnothing$, the strict sublevel is empty, and the exterior
hypotheses for the special closed sublevel cannot hold.
\end{proposition}

\begin{proof}
We first dispose of nonemptiness and the rank-zero convention.  If
$T=\varnothing$, there is no proper subset $J$ in
\eqref{eq:residues-pairs}, so the strict sublevel is empty.  The
closed exterior hypotheses would prescribe a cap and $m^q=c$ on
$I\cup\partial I$.  This is precisely the rank-zero situation
excluded by Proposition~\ref{prop:isolated-layer}.

Now let $T\ne\varnothing$ and assume the closed exterior hypotheses.
Lemma~\ref{lem:cap-accessibility} gives a representative capped on
$T$, and hence on $I$.  If the closed sublevel were empty, then at
every such representative $z$ one would have $m_j^q(z)>c_j$ for
every $j\in T$.  Indeed, any $j$ with $m_j^q(z)\le c_j$ would give
the allowed pair $(T\setminus\{j\},A_{T\setminus\{j\}}z)$.
The resulting strict inequalities at every capped representative
contradict Proposition~\ref{prop:isolated-layer}, applied with its
spherical subset equal to $T$.  Thus the closed sublevel is nonempty.

We prove the contractibility assertions together by induction on
$|T|$, with the conventions just established as the base.  Let
$\mathcal R$ denote either of the two sublevels, and suppose that the
proposition holds for all smaller spherical subsets.  For
$X=(J,A_Jz)\in\mathcal R$, define $p_X$ by
\[
 (p_X)_i=m_i^q(z)\quad(i\notin J),\qquad (Cp_X)_J=0,
 \qquad u_X=(p_X)_T-b_T.
\]
This vector is independent of the representative of $A_Jz$.
On $V=T\setminus J$ it satisfies
$(u_X)_i=m_i^q(z)-b_i\le0$.  Harmonic extension gives
\begin{equation}
 \label{eq:residues-harmonic-difference}
 (u_X)_J=-C_J^{-1}C_{J,V}(u_X)_V\le0,
 \qquad (Cp_X)_T=C_Tu_X.
\end{equation}
The inequality uses nonnegativity of $C_J^{-1}$ and nonpositivity
of the off-diagonal entries of $C$; the formula includes the empty
$J$ convention.  In particular $u_X\le0$ coordinatewise.  In the
strict sublevel $u_X\ne0$, since $V$ is nonempty and all of its
inequalities are strict.

Filter the vertices by increasing height
\[
 h(X)=\sum_{i\in T}(u_X)_i\le0.
\]
Only finitely many height values occur.  For each possible type $J$,
the open counts belong to bounded intervals $0\le m_i^q(z)\le b_i$,
and each such interval contains only finitely many possible layer
counts, by Section~\ref{sec:height-data}.  The exterior vector is
fixed, and harmonic extension determines the remaining coordinates.
There are only finitely many types.  At equal negative height put
larger $|J|$ first; at height zero put smaller $|J|$ first.  Break all
remaining ties by any well-order.  This gives a well-order of
$\mathcal R$.  Remaining ties occur only between incomparable
distinct vertices, so they do not affect any earlier lower or upper
part.

\medskip\noindent
\textit{The upper part at nonzero height.}
Fix $X=(J,A_Jz)$ with $u=u_X\ne0$, and put $V=T\setminus J$.
Every upper extension has the unique form
$(K,A_Kz)$ with $J\subsetneq K\subsetneq T$.  It is allowed because
its open colors are a subset of $V$ and retain the same data.
Writing $p_K$ for its harmonic vector, the change is supported on
$K$ and is
\begin{equation}
 \label{eq:residues-upper-change}
 (p_K-p_X)_K=-C_K^{-1}(C_Tu)_K.
\end{equation}
The rows $(C_Tu)_J$ vanish.  Set
\[
 D=\{i\in V:(C_Tu)_i\ge0\}.
\]
If $K\setminus J$ avoids $D$, all its rows in
\eqref{eq:residues-upper-change} are strictly negative.  The
nonnegative inverse, with positive diagonal, makes the total change
in height strictly positive.  Such an extension is later.  If instead
$\varnothing\ne K\setminus J\subseteq D$, the change is
coordinatewise nonpositive, so the extension is earlier: a strict
height decrease suffices, and equality is resolved in its favor by
the larger type at the negative height $h(X)$.

It follows that every earlier upper extension meets $D$ in a new
color, whereas every extension using only new colors in $D$ is
earlier.  Also $D\ne V$.  Otherwise $C_Tu\ge0$, and
$u=C_T^{-1}C_Tu\ge0$, contradicting $u\le0$ and $u\ne0$.
If $D=\varnothing$, the earlier upper part is empty.  If
$D\ne\varnothing$, define on that part
\[
 (K,A_Kz)\longmapsto
 \bigl(K\cap(J\cup D),A_{K\cap(J\cup D)}z\bigr).
\]
The new type properly contains $J$, uses only new colors in $D$,
and is a proper subset of $T$.  Hence the image is an earlier upper
extension.  This order map is below the identity and retracts onto
the subposet of extensions whose types are contained in $J\cup D$.
That subposet has maximum $(J\cup D,A_{J\cup D}z)$, which is itself
earlier.  Thus the earlier upper part is contractible.

\medskip\noindent
\textit{The lower part at nonzero height.}
Since equal negative heights place larger types first, every earlier
lower pair has strictly smaller height.  Write such a pair as
$Y=(R,A_Rw)$, where $R\subsetneq J$ and $w\in A_Jz$.
Its newly open colors are $U=J\setminus R$.  The vector
$\delta=p_Y-p_X$ vanishes outside $J$, equals
$m_i^q(w)-(p_X)_i$ on $U$, and is harmonic on $R$.  If every one
of these new open differences were nonnegative, nonnegative harmonic
interpolation on $R$ would give $\delta\ge0$ and preclude a strict
height decrease.  Therefore the set
\[
 N(Y)=\{i\in J\setminus R:m_i^q(w)<(p_X)_i\}
\]
is nonempty.  All its colors are capped, since $Y$ is allowed.
Close the other newly open colors by defining
\begin{equation}
 \label{eq:residues-lower-retraction}
 \rho(Y)=\bigl(J\setminus N(Y),A_{J\setminus N(Y)}w\bigr).
\end{equation}
This pair is well-defined on the coset $A_Rw$: both the tests on
$J\setminus R$ and the larger coset are unchanged if $w$ is
replaced by another element of $A_Rw$.

The boundary data on $\partial J$ are capped.  Indeed a neighbor
outside $J$ lies either in $V$ or in $\partial T$, where the cap is
already known.  The harmonic bound for the residue $(J,A_Jz)$ is
exactly $p_X$.  Consequently \eqref{eq:residues-lower-retraction}
belongs to the strict sublevel
$\mathcal R_{<}(J,A_Jz;q)$.  Viewed as a pair inside $A_Tx$, it
is allowed: its new open counts are strictly below $p_X\le b$,
and its old open colors $V$ retain their original inequalities.
Its harmonic difference from $p_X$ is strictly negative on the
nonempty set $N(Y)$ and is a nonnegative interpolation of these
differences on the remaining colors of $J$.  It is therefore
coordinatewise nonpositive with strictly negative sum.  Thus
$\rho(Y)$ remains an earlier lower pair.

For completeness, this closure respects actual coset inclusions
when representatives differ.  Suppose
\[
 (R_1,A_{R_1}w_1)\le(R_2,A_{R_2}w_2).
\]
Then $R_1\subseteq R_2$ and $w_1\in A_{R_2}w_2$.  On their
shared open colors $J\setminus R_2$ the counts agree, whence
\[
 N(Y_2)=N(Y_1)\cap(J\setminus R_2)\subseteq N(Y_1).
\]
Writing $K_a=J\setminus N(Y_a)$, we obtain $K_1\subseteq K_2$.
Moreover $R_2\subseteq K_2$, so $w_1\in A_{K_2}w_2$ and
\[
 A_{K_1}w_1\subseteq A_{K_2}w_2.
\]
This proves that $\rho$ is an order map.  It lies above the identity
because $R\subseteq J\setminus N(Y)$.

Conversely, any pair in $\mathcal R_{<}(J,A_Jz;q)$ is an allowed
earlier lower pair in $\mathcal R$: the same interpolation argument
gives a strict decrease of height, and $p_X\le b$ gives the outer
inequalities.  All its newly open colors pass the strict test, so
$\rho$ fixes it.  Thus $\rho$ is an order retraction onto exactly
that strict sublevel.  Since $J\subsetneq T$, induction says that
the earlier lower part is empty or contractible.  This includes
$J=\varnothing$, when both lower posets are empty.

\medskip\noindent
\textit{The zero-height case.}
If $u_X=0$, the pair lies in the special closed sublevel.  All open
colors have $m_i^q(z)=c_i$ and are capped.  An upper extension retains
$p_X=b$, so its height remains zero and its larger type makes it
later.  There is no earlier upper part.  Every allowed lower pair is
earlier: its height is either negative, or zero with a smaller type.

The lower part is exactly
$\mathcal R_{\le}(J,A_Jz;q,I,c)$.  Its exterior hypotheses follow
from the original ones together with the equalities and caps on
$V=T\setminus J$: these sets supply all of
$(I\setminus J)\cup\partial I$.  If $J=\varnothing$, those
hypotheses contradict the rank-zero case.  Hence $J\ne\varnothing$,
and induction makes this earlier lower part nonempty and
contractible.

\medskip\noindent
\textit{Sources.}
We have shown that every earlier lower and upper part is empty or
contractible.  Lemma~\ref{lem:filtration-principle} therefore makes
each component contractible with exactly one source.  The zero-height
case has a nonempty earlier lower part, so a source must have
$u_X\ne0$.  The upper calculation then forces $D=\varnothing$,
and the strict lower sublevel for $(J,A_Jz)$ must be empty.

These facts give all the hypotheses of a candidate in
Proposition~\ref{prop:candidate-uniqueness}, for the fixed ambient
coset $A_Tx$ and cap $q$.  The colors $V$ are capped by definition.
Equation~\eqref{eq:residues-harmonic-difference} gives
\[
 -(Cp_X)_i=-(C_Tu_X)_i>0\quad(i\in V),
 \qquad (Cp_X)_J=0.
\]
Finally consider \emph{any} representative $w\in A_Jz$ capped on
$T$.  If $m_j^q(w)<(p_X)_j$ for some $j\in J$, the pair
$(J\setminus\{j\},A_{J\setminus\{j\}}w)$ would belong to the
strict lower sublevel.  Its emptiness therefore proves
\[
 m_j^q(w)\ge(p_X)_j\qquad(j\in J)
\]
at every such capped representative, as required.  For
$J=\varnothing$ this condition is vacuous.

Proposition~\ref{prop:candidate-uniqueness} gives at most one source
in $\mathcal R$.  A nonempty $\mathcal R$ has a first vertex in
the filtration and hence a source.  It therefore has one component
and is contractible.  The special closed sublevel was already shown
to be nonempty.  This completes both inductions.
\end{proof}
